% Part: first-order-logic
% Chapter: completeness
% Section: maximally-consistent-sets

% Definition of maximally consistent sets. Properties of mcs required
% for completeness proved are in maximally-consistent-sets.tex in the
% chapter on the proof system used.

\documentclass[../../../include/open-logic-section]{subfiles}

\begin{document}

\olfileid[pt]{fol}{com}{mcs}
\olsection{Conjuntos Maximamente Consistentes de \usetoken{P}{sentence}}

\begin{defn}[Conjunto maximamente consistente]
\ollabel{mcs}
Um conjunto~$\Gamma$ de !!{sentence}s é \emph{maximamente consistente} se e somente se
\begin{enumerate}
\item $\Gamma$ é consistente, e
\item se $\Gamma \subsetneq \Gamma'$, então $\Gamma'$ é inconsistente.
\end{enumerate}
\end{defn}

\begin{explain}
Uma definição alternativa, equivalente à anterior, é: um conjunto $\Gamma$ de
sentenças é \emph{maximamente consistente} se e somente se
\begin{enumerate}
\item $\Gamma$ é consistente, e
\item Se $\Gamma \cup \{ !A \}$ é consistente, então $!A \in \Gamma$.
\end{enumerate}
Em outras palavras, não se pode adicionar !!{sentence}s que ainda não estejam em~$\Gamma$
a um conjunto maximamente consistente~$\Gamma$ sem tornar inconsistente o
conjunto maior resultante.
\end{explain}

\begin{explain}
Conjuntos maximamente consistentes são importantes na prova de completude
já que podemos garantir que todo conjunto consistente de !!{sentence}s~$\Gamma$
está contido em um conjunto maximamente consistente~$\Gamma^*$, e um conjunto maximamente
consistente contém, para cada !!{sentence}~$!A$, ou $!A$ ou sua
negação $\lnot !A$. Isso é verdade em particular para !!{sentence}s atômicas,
de modo que, a partir de um conjunto maximamente consistente em uma linguagem adequadamente expandida por
!!{constant}s, podemos construir !!a{structure} onde a
interpretação dos !!{predicate}s é definida de acordo com quais !!{sentence}s
atômicas estão em~$\Gamma^*$. Pode-se então mostrar que esta !!{structure}
torna verdadeiras todas as !!{sentence}s em~$\Gamma^*$ (e, portanto, também em $\Gamma$).
A prova deste último fato requer que $\lnot !A \in
\Gamma^*$ se e somente se $!A \notin \Gamma^*$, $(!A \lor !B) \in \Gamma^*$ se e somente se $!A
\in \Gamma^*$ ou $!B \in \Gamma^*$, etc.
\end{explain}

\begin{prop}
\ollabel{prop:mcs}
Suponha que $\Gamma$ seja maximamente consistente. Então:
\begin{enumerate}
\item \ollabel{prop:mcs-prov-in} Se $\Gamma \Proves !A$, então $!A \in
  \Gamma$.

\item \ollabel{prop:mcs-either-or} Para qualquer $!A$, ou $!A \in
  \Gamma$ ou $\lnot !A \in \Gamma$.

\tagitem{prvAnd}{\ollabel{prop:mcs-and} $(!A \land !B) \in \Gamma$
  se e somente se $!A \in \Gamma$ e $!B \in \Gamma$.}{}

\tagitem{prvOr}{\ollabel{prop:mcs-or} $(!A \lor !B) \in \Gamma$ se e somente se
  $!A \in \Gamma$ ou $!B \in \Gamma$.}{}

\tagitem{prvIf}{\ollabel{prop:mcs-if} $(!A \lif !B) \in \Gamma$ se e somente se
  $!A \notin \Gamma$ ou $!B \in \Gamma$.}{}
\end{enumerate}
\end{prop}

\begin{proof}
Suponhamos, para tudo o que se segue, que $\Gamma$ seja maximamente
consistente.
\begin{enumerate}
\item Se $\Gamma \Proves !A$, então $!A \in \Gamma$.

Suponha que $\Gamma \Proves !A$. Suponha, ao contrário, que $!A
\notin \Gamma$: então, como $\Gamma$ é maximamente consistente, $\Gamma
\cup \{!A\}$ é inconsistente, logo $\Gamma \cup \{!A\} \Proves
\lfalse$. Por
\tagrefs{prfSC/{fol:seq:prv:prop:provability-contr},prfND/{fol:ntd:prv:prop:provability-contr}},
$\Gamma$ é inconsistente. Isso contradiz a suposição de que
$\Gamma$ é consistente. Logo, não pode ser o caso que $!A \notin
\Gamma$, de modo que $!A \in \Gamma$.

\item Para qualquer $!A$, ou $!A \in \Gamma$ ou $\lnot !A \in \Gamma$.

Suponha, ao contrário, que para algum~$!A$ ambos $!A \notin \Gamma$ e
$\lnot !A \notin \Gamma$. Como $\Gamma$ é maximamente consistente,
$\Gamma \cup \{!A\}$ e $\Gamma \cup \{\lnot !A\}$ são ambos
inconsistentes, logo $\Gamma \cup \{!A\} \Proves \lfalse$ e $\Gamma \cup
\{\lnot !A\} \Proves \lfalse$. Por
\tagrefs{prfSC/{fol:seq:prv:prop:provability-exhaustive},prfND/{fol:ntd:prv:prop:provability-exhaustive}},
$\Gamma$ é inconsistente, uma contradição. Logo, não pode haver tal
!!a{sentence}~$!A$ e, para todo $!A$, $!A \in \Gamma$ ou $\lnot !A
\in \Gamma$.

\tagitem{defAnd}{}{%
\iftag{probAnd}{Exercício.}{%
$(!A \land !B) \in \Gamma$ se e somente se $!A \in \Gamma$ e $!B \in \Gamma$:

Para a direção direta, suponha $(!A \land !B) \in \Gamma$. Então
$\Gamma \Proves !A \land !B$. Por
\tagrefs{prfSC/{fol:seq:prv:prop:provability-land-left},prfND/{fol:ntd:prv:prop:provability-land-left}},
$\Gamma \Proves !A$ e $\Gamma \Proves !B$. Por
\olref{prop:mcs-prov-in}, $!A \in \Gamma$ e $!B \in \Gamma$, como
requerido.

Para a direção reversa, sejam $!A \in \Gamma$ e $!B \in
\Gamma$. Então $\Gamma \Proves !A$ e $\Gamma \Proves !B$. Por
\tagrefs{prfSC/{fol:seq:prv:prop:provability-land-right},prfND/{fol:ntd:prv:prop:provability-land-right}},
$\Gamma \Proves !A \land !B$. Por \olref{prop:mcs-prov-in}, $(!A \land
!B) \in \Gamma$.}}

\tagitem{defOr}{}{%
\iftag{probOr}{Exercício.}{%
$(!A \lor !B) \in \Gamma$ se e somente se $!A \in \Gamma$ ou $!B \in \Gamma$.

Para a contrapositiva da direção direta, suponha que $!A
\notin \Gamma$ e $!B \notin \Gamma$. Queremos mostrar que $(!A \lor
!B) \notin \Gamma$. Como $\Gamma$ é maximamente consistente, $\Gamma
\cup \{!A\} \Proves \lfalse$ e $\Gamma \cup \{!B\} \Proves \lfalse$.
Por
\tagrefs{prfSC/{fol:seq:prv:prop:provability-lor-left},prfND/{fol:ntd:prv:prop:provability-lor-left}},
$\Gamma \cup \{(!A \lor !B)\}$ é inconsistente. Logo, $(!A \lor !B)
\notin \Gamma$, como requerido.

Para a direção reversa, suponha que $!A \in \Gamma$ ou $!B \in
\Gamma$. Então $\Gamma \Proves !A$ ou $\Gamma \Proves !B$. Por
\tagrefs{prfSC/{fol:seq:prv:prop:provability-lor-right},prfND/{fol:ntd:prv:prop:provability-lor-right}},
$\Gamma \Proves !A \lor !B$. Por \olref{prop:mcs-prov-in}, $(!A \lor
!B) \in \Gamma$, como requerido.}}

\tagitem{defIf}{}{%
\iftag{probIf}{Exercício.}{%
$(!A \lif !B) \in \Gamma$ se e somente se $!A \notin \Gamma$ ou $!B \in \Gamma$:

Para a direção direta, seja $(!A \lif !B) \in \Gamma$ e suponha,
ao contrário, que $!A \in \Gamma$ e $!B \notin \Gamma$. Sob essas
suposições, $\Gamma \Proves !A \lif !B$ e $\Gamma \Proves !A$. Por
\tagrefs{prfSC/{fol:seq:prv:prop:provability-mp},prfND/{fol:ntd:prv:prop:provability-mp}},
$\Gamma \Proves !B$. Mas então, por \olref{prop:mcs-prov-in}, $!B \in
\Gamma$, contradizendo a suposição de que $!B \notin \Gamma$.

Para a direção reversa, considere primeiro o caso em que $!A \notin
\Gamma$. Por \olref{prop:mcs-either-or}, $\lnot !A \in \Gamma$ e
portanto $\Gamma \Proves \lnot !A$. Por
\tagrefs{prfSC/{fol:seq:prv:prop:provability-lif},prfND/{fol:ntd:prv:prop:provability-lif}},
$\Gamma \Proves !A \lif !B$. Novamente por \olref{prop:mcs-prov-in}, obtemos
que $(!A \lif !B) \in \Gamma$, como requerido.

Agora considere o caso em que $!B \in \Gamma$. Então $\Gamma \Proves !B$
e por
\tagrefs{prfSC/{fol:seq:prv:prop:provability-lif},prfND/{fol:ntd:prv:prop:provability-lif}},
$\Gamma \Proves !A \lif !B$. Por \olref{prop:mcs-prov-in}, $(!A \lif
!B) \in \Gamma$.}}
\end{enumerate}
\end{proof}

\begin{prob}
Complete a prova de \olref[fol][com][mcs]{prop:mcs}.
\end{prob}

\end{document}
